\documentclass[../../../thesis.tex]{subfiles}
\begin{document}
\subsection{Matematické rovnice}
Systém na sadzbu textu \TeX\ pôvodne vyvinul Donald Knuth.
Jeho motivácia bola poskytnúť producentom vedeckej tlače
počítačový nástroj,
ktorý bude správne sádzať matematické rovnice.
\LaTeX, ako jeho nadstavba, má sadzbu rovníc takpovediac
v~genetickej výbave.
Autori textov z prírodovedeckej a technickej komunity siahajú po tomto nástroji
práve z~dôvodu bezkonkurenčnej práce s~rovnicami pri tvorbe
vedeckého alebo akademického obsahu.

Matematické rovnice používame v tlačenom texte dvomi spôsobmi:
1. píšeme ich v~rámci textového odseku;
2. rovnicu vytlačíme zvlášť medzi dva textové odseky
a~vtedy ju spravidla aj číslujeme, aby sme sa na ňu mohli
ďalej odvolávať.

\subsubsection{Rovnica v textovom riadku}
Riešenie kvadratickej rovnice s koeficientami $a,b,c$
a~s~neznámou $x$ vypočítame pomocou známeho vzťahu
$x=\frac{-b\pm\sqrt{b^2-4ac}}{2a}$.
Je to príklad rovnice zapísanej v~rámci textového odseku.
Ak tú istú rovnicu napíšeme do samostatného odseku,
vyzerá trochu inak:
$$
  x=\frac{-b\pm\sqrt{b^2-4ac}}{2a}
$$

Očividný rozdiel je vo veľkosti zlomku a~znaku odmocniny,
môžeme si všimnúť aj malé rozdiely v~medzerách,
vo vertikálnom zarovnávaní, atď.

Vložené rovnice v~rámci textového riadku zapisujeme pomocou
znaku \$.
Matematický zápis ohraničíme znakmi dolára sprava aj zľava.
Napríklad zápis \verb|$y = ax^2 + bx + c$| vytvorí rovnicu $y = ax^2 + bx + c$. Rovnakú funkciu ako znak dolára má dvojica znakov \verb|\(...\)|.

Označenia fyzikálnych veličín píšeme tiež ako vloženú rovnicu:
veľkosť sily $F$, hmotnosť $m$, čas $t$ a~podobne.
Všetky veličiny sme zapísali takto: \verb|$F$|,
\verb|$m$|, \verb|$t$|.

\subsubsection{Zobrazená rovnica}
Matematický text ohraničený dvomi znakmi dolára zľava
a~dvomi dolármi sprava interpretuje kompilátor \LaTeX-u
ako zobrazenú rovnicu, ktorú vysádza do zvláštneho odseku
zarovnaného na stred, napríklad
$$
  y = ax^2 + bx + c
$$
sme získali kódom \verb|$$y = ax^2 + bx + c$$|.
Namiesto dolárov môžeme použiť dvojicu príkazov
\verb|\[...\]| s rovnakým výsledkom.

Rovnicu s referenčným číslom vytvoríme tak, že zápis matematickej
rovnice vložíme do prostredia \verb|equation|:
\begin{verbatim}
  \begin{equation}
    y = ax^2 + bx + c
  \end{equation}
\end{verbatim}
vytvorí rovnicu s číslom uzavretým v zátvorkách:
\begin{equation}
  y = ax^2 + bx + c
\end{equation}

Program \LaTeX\ čísluje rovnice automaticky od čísla 1.

\subsubsection{Zásady matematickej sadzby}
Pravidlá sadzby matematických, fyzikálnych veličín a~ich vzťahov
sumarizuje medzinárodná norma u~nás známa pod označením
STN ISO 80 000: 2022 Veličiny a~jednotky \cite{iso800001, iso800002}.
Označenie fyzikálnych a~matematických veličín píšeme vždy
šikmým rezom písma.
Čísla, názvy funkcií a~jednotky fyzikálnych veličín zapisujeme
normálnym rezom.
Správny zápis elektrického napätia s veľkosťou 5,07 voltu
vyzerá takto:
\begin{equation}\label{Eq:quantity}
  U = 5{,}07\,\mathrm{V}
\end{equation}
kde $U$ je elektrické napätie.
Môžeme si všimnúť, že okolo znaku rovnosti sú medzery,
desatinná čiarka sa píše bez medzier a~medzi číslom a~jednotkou
je úzka medzera -- v~\LaTeX-u príkaz \verb|\,|.

Rovnicu (\ref{Eq:quantity}) sme zapísali v zdrojovom kóde nasledujúcim spôsobom:
\begin{verbatim}
\begin{equation}
  U = 5{,}07\,\mathrm{V}
\end{equation}
\end{verbatim}
\TeX\ v~matematickom móde automaticky sádže veličiny kurzívou.
Ak chceme, aby bola jednotka V vzpriamená, použijeme
v~matematickom móde príkaz \verb|\mathrm{}|. Medzery okolo znaku rovnosti sú taktiež automatické.
Zaujímavá je desatinná čiarka, ktorú musíme uzavrieť do
zložených zátvoriek, aby sme potlačili automatickú sadzbu
medzery za čiarkou.
V~anglicky hovoriacich krajinách sa ako desatinný oddeľovač
používa bodka.
Čiarka má väčšinou význam oddeľovača prvkov zoznamov
a~v~matematickom režime vkladá \TeX\, za čiarku úzku medzeru.
Najčastejšie chyby pri zápise fyzikálnych veličín sme zhrnuli
v~tabuľke 1.

\begin{table}[h]
\centering
\caption{Prehľad najčastejších chýb pri nesprávnom zápise
skalárnej fyzikálnej veličiny.
Správny zápis predstavuje rovnica (\ref{Eq:quantity}).}
\vspace{12pt}
\begin{tabular}{@{}ll@{}}
  \toprule
  \textbf{nesprávny zápis} & \textbf{opis chyby}\\
  \midrule
  $U = 5{,}07\,V$ & jednotka je kurzívou\\
  $U = 5{,}07\mathrm V$ & medzi číslom a jednotkou chýba medzera\\
  $\mathrm U = 5{,}07\,\mathrm V$ & označenie veličiny nie je kurzívou \\
  $U$=$5{,}07\,\mathrm V$ & okolo znaku rovnosti chýbajú medzery\\
  $U = 5,07\,\mathrm{V}$ & za desatinnou čiarkou je medzera\\
  $U = \mathit{5{,}07}\,\mathrm{V}$ & číslo je kurzívou\\
  U=$\mathit{5,07}V$ & kumulácia predchádzajúcich chýb\\
  \bottomrule
\end{tabular}
\end{table}

\paragraph{Dôležité pravidlá písania rovníc}
\begin{itemize}

    \item Značky veličín píšeme šikmým rezom písma (kurzívou):
    $x$, $y$, $a$, $F$, $P$, $W$.

    \item Fyzikálne jednotky píšeme vzpriameným písmom:
    $a = 10\,\mathrm{cm}$.

    \item Čísla píšeme vzpriameným písmom:
    $1$; $2$; $3$; $1\,024$; $3{,}14$ a podobne.

    \item Skratky matematických funkcií píšeme vzpriameným
    písmom: $\sin(\alpha + \beta)$, $\cos\omega t$,
    $\log_a x = \frac{\ln x}{\ln a}$, $\mathrm{e^{i\uppi}}=-1$.

    \item Označenia nemenných konštánt sú tiež vzpriamené
    písmená: $\uppi$, $\mathrm i$, $\mathrm e$ --
    tri základné matematické konštanty --
    Ludolfovo číslo, komplexná jednotka a~Eulerovo číslo.
    Niektoré konštanty sa zo zvyku môžu písať kurzívou,
    napríklad $\pi$ alebo dielektrická konštanta $\varepsilon_0$.
    Komplexná jednotka je však vždy vzpriamená:
    $\mathrm i^2 = -1$.

    \item Vzpriameným písmom píšeme v matematických vzťahoch
    aj všetky zátvorky.

    \item Sumačné indexy píšeme kurzívou:
    $p_N(x) = \sum_{i=1}^N a_ix^i$.
    Symbol $i$ v tomto príklade predstavuje sumačný index,
    nie komplexnú jednotku.

    \item Vektory uvádzame buď polotučným šikmým rezom
    ($\vectorsym a$, $\vectorsym b$, $\vectorsym F$)
    alebo šikmým netučným rezom so šípkou nad symbolom:
    $\vec a$, $\vec b$, $\vec F$.
    Treba si vybrať jeden spôsob a~ten používať v~celej práci.

    Označenia matíc a tenzorov zapisujeme polotučným šikmým rezom:
    $$\matrixsym M = \left( \begin{array}{cc}
                            m_{11} & m_{12}\\
                            m_{21} & m_{22}
                        \end{array}\right)
    $$
    Prvky matice $m_{ij}$ sú skalárne veličiny, preto sú to
    netučné šikmé písmená.

    \item Ak treba z nejakého dôvodu odlíšiť tenzor
    od bežnej matice, môžeme tenzory označiť dvomi čiarkami:
    $\overline{\overline T}$.

    \item Značku úplného diferenciálu píšeme vzpriameným rezom:
    $\mathrm dy$ je úplný diferenciál veličiny $y$.

    \item Derivácia dráhy podľa času:
    $v = \frac{\mathrm ds}{\mathrm dt}$.
    Veličiny $v$, $s$ a~$t$ sú stále písané kurzívou

    \item Určitý integrál vyzerá takto:
    $$
    \int_a^b f(x)\,\mathrm dx
    $$
    V~integráli spravidla vkladáme pred diferenciál
    úzku medzeru \verb|\,|.
\end{itemize}

\subsubsection*{\normalsize Príklad}
Z Coulombovho zákona vyplýva,
že pre vektor elektrostatickej sily $\mathbfit F_\mathrm e$
medzi dvomi bodovými nábojmi platí nasledujúci vzťah
\begin{equation}\label{Eq:CoulombVec}
  \mathbfit F_\mathrm e = \frac1{4\pi\varepsilon_0}
  \frac{q_1q_2}{r^2}\frac{\mathbfit r}r
\end{equation}
kde $q_1$, $q_2$ sú veľkosti bodových nábojov,
$\mathbf r$ je polohový vektor náboja $q_2$
vzhľadom na náboj $q_1$
a~$\varepsilon_0$ je elektrická konštanta.

Aby sme zhrnuli predchádzajúce pravidlá, detailnejšie opíšeme spôsob zápisu jednotlivých prvkov
v~rovnici (\ref{Eq:CoulombVec}).
Skalárne veličiny veľkosť náboja a~vzájomná vzdialenosť sú
napísané kurzívou,
vektorové veličiny sila a~polohový vektor sú polotučným rezom.
Všetky čísla (indexy a~násobok 4 v~menovateli)
píšeme normálnym rezom.
Konštanty $\pi$ a $\varepsilon_0$ sú podľa zvyklosti vysádzané
šikmým rezom.
Norma v~súčasnosti odporúča vzpriamené písmo:
$\uppi$, $\upvarepsilon_0$, čo možno dosiahnuť použitím balíka
\verb|upgreek|.
Index $\mathrm e$ pri symbole vektora sily,
ktorý označuje fakt, že ide o~elektrickú silu,
píšeme normálnym neskloneným rezom písma.

V~texte, ktorý nasleduje bezprostredne za rovnicou vysvetlíme
a~stručne opíšeme jednotlivé symboly.
Tento odsek formálne patrí k~rovnici,
preto nemá odsadený prvý riadok.
Ak to chceme dosiahnuť,
nevynecháme za rovnicou v~zdrojovom kóde
prázdny riadok.

Zdrojový kód rovnice \eqref{Eq:CoulombVec}:
\begin{verbatim}
\begin{equation}\label{Eq:CoulombVec}
  \mathbfit F_\mathrm e = \frac1{4\pi\varepsilon_0}
  \frac{q_1q_2}{r^2}\frac{\mathbfit r}r
\end{equation}
\end{verbatim}
Všimnite si, že niekedy argument funkcie nemusí byť
uzavretý zloženými zátvorkami (\verb|\mathrm e| je vo výsledku rovnaký ako \verb|\mathrm{e}|).
Tento fenomén je dôsledok priebehu kompilácie na úrovni tzv. procesorov \TeX-u.
Parameter bez zátvorky funguje správne vtedy,
ak je parameter práve jeden \emph{token}.%
\footnote{Čo je token? Je to akási najmenšia časť spracovaného
kódu, ktorý vstupuje do tzv. \emph{token procesora}
počas kompilácie.
Tokeny sú písmená, čísla, ale aj jednoduchý príkaz,
napr. \texttt{\textbackslash mathrm} je jeden token.
Problematika je síce komplexná, ale na pochopenie
toho, čo sa v~\TeX-u deje, dôležitá.
Existuje veľmi pekná kniha od Petra Olšáka \TeX book naruby z roku 2001 \cite{Olsak2001Texbook}.
Keďže sa už nedá zohnať, autor sa rozhodol, že ju sprístupní pre širokú verejnosť v digitálnej forme na stránke
\href{https://petr.olsak.net/tbn.html}{\texttt{petr.olsak.net/tbn.html}}.
}
\end{document}